\documentclass[12pt]{article}

\usepackage{hcmus-report-template}

% Disable indentation on new paragraphs
\setlength{\parindent}{0pt}

% Line spacing 1.5
\renewcommand{\baselinestretch}{1.5}

% Graphic path
\graphicspath{{img/}}

% Define course name, report name and report title.
\newcommand{\coursename}{DevOps}
\newcommand{\reportname}{Triển khai CI/CD, GitOps, Service Mesh và Observability cho hệ thống YAS --- Yet Another Shop}
\newcommand{\reporttitle}{Báo Cáo Đồ Án 2: CI/CD, GitOps, Service Mesh và Observability}

\newcommand{\studentname}{%
  Ngô Gia An (23120205)\\%
  Nguyễn Bảo An (23120207)\\%
  Dương Tuấn Anh (23120208)\\%
  Trương Nhật Đạt (23120231)%
}
\newcommand{\teachername}{Giáo viên hướng dẫn}

% Header
\lhead{\reporttitle}
\rhead{%
  Trường Đại học Khoa học Tự nhiên - ĐHQG HCM\\%
  \coursename%
}

% Footer
\newcommand{\leftfooter}{Nhóm: com-suon-bi-cha}
\lfoot{\leftfooter}

% ============ DOCUMENT ============
\begin{document}

\pagenumbering{roman}
\input{content/title.tex}

\tableofcontents
\pagebreak
\listoffigures
\pagebreak

\pagenumbering{arabic}
\setcounter{page}{1}

\section{Thông Tin Nhóm}

\begin{table}[H]
\centering
\caption{Thông tin thành viên nhóm}
\begin{tabular}{|l|c|p{0.52\textwidth}|}
\hline
\textbf{Thành viên} & \textbf{MSSV} & \textbf{Phần phụ trách} \\
\hline
Ngô Gia An & 23120205 & Hạ tầng GCP/K3s, ArgoCD, Jenkins agent và Observability \\
\hline
Nguyễn Bảo An & 23120207 & Jenkins CI/CD, build/push image và developer-build \\
\hline
Trương Nhật Đạt & 23120231 & GitOps manifest, Kustomize overlays và ArgoCD applications \\
\hline
Dương Tuấn Anh & 23120208 & Istio Service Mesh, security policy, Kiali và minh chứng vận hành \\
\hline
\end{tabular}
\end{table}

\textbf{Link GitHub Repository:} \url{https://github.com/com-suon-bi-cha/yas}

\textbf{Repo GitOps manifest:} \url{https://github.com/com-suon-bi-cha/gitops-manifest-k8s}

\section{Yêu Cầu Đồ Án}

Phần này tổng hợp các yêu cầu chính của đồ án và trạng thái hoàn thành tương ứng. Các nội dung được đối chiếu theo kết quả triển khai thực tế trên Jenkins, repo GitOps, ArgoCD và Kubernetes cluster.

\begin{table}[H]
\centering
\caption{Đối chiếu yêu cầu chính của đồ án}
\small
\begin{tabular}{|p{0.30\textwidth}|p{0.46\textwidth}|c|}
\hline
\textbf{Yêu cầu} & \textbf{Kết quả thực hiện} & \textbf{Trạng thái} \\
\hline
CI cho các service trong scope & Jenkins build/test/push Docker image cho các workload ứng dụng được chọn & Đạt \\
\hline
CD lên Kubernetes & ArgoCD đồng bộ manifest từ repo GitOps vào namespace \texttt{dev} và \texttt{staging} & Đạt \\
\hline
GitOps manifest & Kustomize quản lý base và overlays cho \texttt{dev}, \texttt{staging}, \texttt{developer-build} & Đạt \\
\hline
Service Mesh nâng cao & Istio cấu hình mTLS, retry, AuthorizationPolicy, ingress routing và Kiali visualization cho môi trường \texttt{dev}; \texttt{staging} có overlay Istio riêng & Đạt \\
\hline
Observability & Prometheus, Grafana, Loki, Tempo, Promtail và OpenTelemetry Collector được triển khai trong namespace \texttt{observability} & Đạt \\
\hline
Minh chứng vận hành & Screenshot Jenkins, ArgoCD, Kubernetes, Kiali và Observability được tổ chức trong thư mục báo cáo & Đạt \\
\hline
\end{tabular}
\end{table}

Phạm vi triển khai hiện tại gồm 16 workload ứng dụng:

\begin{itemize}
  \item Core services: \texttt{product}, \texttt{cart}, \texttt{order}, \texttt{customer}, \texttt{inventory}, \texttt{tax}, \texttt{payment}.
  \item Supporting services: \texttt{media}, \texttt{search}, \texttt{location}, \texttt{sampledata}.
  \item Storefront: \texttt{storefront-bff}, \texttt{storefront-ui}.
  \item Backoffice và công cụ: \texttt{backoffice-bff}, \texttt{backoffice-ui}, \texttt{swagger-ui}.
\end{itemize}

Workload \texttt{kafka-connect} được triển khai như thành phần hỗ trợ CDC, không tính là service ứng dụng chính.

\section{Phân Công Công Việc}

Việc phân công được tổ chức theo các lớp kỹ thuật của hệ thống. Mỗi phạm vi có đầu ra độc lập, đồng thời phụ thuộc vào các phần còn lại để hình thành chuỗi CI/CD hoàn chỉnh.

\begin{table}[H]
\centering
\caption{Phân công theo phạm vi kỹ thuật}
\begin{tabular}{|p{0.2\textwidth}|p{0.32\textwidth}|p{0.34\textwidth}|}
\hline
\textbf{Phụ trách} & \textbf{Phạm vi kỹ thuật} & \textbf{Sản phẩm chính} \\
\hline
Ngô Gia An & Hạ tầng cloud, K3s, namespace, secret, ArgoCD, Jenkins agent và Observability & Cluster sẵn sàng cho CI/CD, GitOps, Service Mesh và hệ thống quan sát \\
\hline
Nguyễn Bảo An & Jenkins CI/CD, Docker image, developer-build và cập nhật image tag & Pipeline build/test/push và cơ chế triển khai theo branch \\
\hline
Trương Nhật Đạt & GitOps manifest, Kustomize, overlay môi trường và ArgoCD application & Repo manifest có thể đồng bộ vào \texttt{dev}, \texttt{staging}, \texttt{developer-build} \\
\hline
Dương Tuấn Anh & Istio Service Mesh, mTLS, AuthorizationPolicy, retry, Kiali và ingress verification & Cấu hình mesh cho \texttt{dev}, kiểm chứng allow/deny, topology và ingress \\
\hline
\end{tabular}
\end{table}

\section{Tổng Quan Kiến Trúc Pipeline}

Hệ thống YAS được triển khai theo mô hình CI/CD kết hợp GitOps. Source code ứng dụng và manifest triển khai được tách thành hai nguồn quản lý độc lập. Khi có thay đổi mã nguồn, Jenkins Multibranch Pipeline thực hiện build/test, đóng gói Docker image và push image lên Docker Hub. Sau đó pipeline cập nhật image tag trong repo \texttt{gitops-manifest-k8s}. ArgoCD theo dõi repo GitOps và đồng bộ trạng thái mong muốn vào Kubernetes cluster.

\begin{figure}[H]
\centering
\fbox{\begin{minipage}{0.88\textwidth}
\centering
\texttt{GitHub} $\rightarrow$ \texttt{Jenkins CI/CD} $\rightarrow$ \texttt{Docker Hub} $\rightarrow$ \texttt{GitOps Manifest} $\rightarrow$ \texttt{ArgoCD} $\rightarrow$ \texttt{K3s Cluster}
\end{minipage}}
\caption{Luồng triển khai tổng quát của Project 2}
\end{figure}

Trong cluster, namespace \texttt{dev} và \texttt{staging} đều được quản lý bằng GitOps. Namespace \texttt{dev} là môi trường kiểm thử chính cho Service Mesh, còn \texttt{staging} có overlay Istio riêng để đáp ứng yêu cầu triển khai chính sách mesh ở môi trường tiền triển khai. Hai môi trường này có sidecar injection, mTLS, retry policy và AuthorizationPolicy. Namespace \texttt{developer-build} phục vụ triển khai tạm thời cho branch khi chạy job kiểm thử. Namespace \texttt{observability} chứa các thành phần quan sát hệ thống như Prometheus, Grafana, Loki, Tempo, Promtail và OpenTelemetry Collector.

\section{Hạ Tầng Kubernetes Và GitOps}

\subsection{Nền tảng GCP và K3s}

Hạ tầng triển khai của đồ án được đặt trên một máy ảo Google Cloud Platform. Máy ảo này chạy K3s và đóng vai trò control-plane đồng thời schedule workload ứng dụng trong giai đoạn demo. Mô hình này phù hợp với phạm vi đồ án vì giảm độ phức tạp vận hành, vẫn đảm bảo đủ môi trường để kiểm chứng CI/CD, GitOps, Service Mesh và Observability. Jenkins agent chạy trên máy local của nhóm và sử dụng kubeconfig để truy cập cluster khi pipeline cần build, cập nhật manifest hoặc kiểm thử luồng developer-build.

\begin{figure}[H]
\centering
\reportimg{member1-report/01-gcp-vm-running.jpg}
\caption{Máy ảo GCP đang ở trạng thái Running và cung cấp địa chỉ public để truy cập cluster}
\end{figure}

Các firewall rule được cấu hình để mở các cổng cần thiết cho SSH, Kubernetes API và NodePort/ingress phục vụ truy cập ứng dụng, ArgoCD, Grafana và các công cụ kiểm chứng. Các cổng này chỉ phục vụ mục tiêu triển khai và demo trong phạm vi đồ án.

\begin{figure}[H]
\centering
\reportimg{member1-report/02-gcp-firewall-rules.jpg}
\caption{Firewall rules trên GCP mở các cổng cần thiết cho SSH, Kubernetes API, NodePort và dịch vụ quan sát}
\end{figure}

\begin{figure}[H]
\centering
\reportimg{member1-report/03-k3s-node-ready.jpg}
\caption{Node K3s ở trạng thái Ready, chứng minh cluster có thể schedule workload}
\end{figure}

\subsection{Namespace và Secret dùng chung}

Cluster được chia theo namespace để tách vai trò vận hành. Namespace \texttt{dev} là môi trường kiểm thử chính và được ArgoCD đồng bộ liên tục. Namespace \texttt{staging} dùng cho kiểm thử release với hạ tầng riêng. Namespace \texttt{developer-build} phục vụ triển khai tạm thời sau khi Jenkins build branch phát triển; môi trường này dùng chung hạ tầng với \texttt{dev} và có thể được cleanup sau khi kiểm thử. Các namespace \texttt{argocd}, \texttt{istio-system} và \texttt{observability} lần lượt chứa GitOps controller, Service Mesh và hệ thống quan sát.

\begin{figure}[H]
\centering
\reportimg{member1-report/04-kubernetes-namespaces.jpg}
\caption{Các namespace chính của đồ án gồm dev, staging, developer-build, argocd, istio-system và observability}
\end{figure}

Docker registry credential được tạo dưới dạng Kubernetes secret trong các namespace triển khai. Secret này cho phép Pod pull image từ Docker Hub trong quá trình rollout, đồng thời tránh lưu thông tin đăng nhập trực tiếp trong manifest GitOps.

\begin{figure}[H]
\centering
\reportimg{member1-report/05-dockerhub-pull-secrets.jpg}
\caption{Docker Hub pull secret tồn tại trong các namespace dev, staging và developer-build}
\end{figure}

\subsection{ArgoCD và trạng thái GitOps}

ArgoCD được triển khai trong namespace \texttt{argocd} để duy trì desired state của các môi trường GitOps. Các thành phần chính như application controller, repo server, API server, Dex, Redis và notifications controller đều chạy ổn định. Khi manifest trong repo GitOps thay đổi, ArgoCD phát hiện khác biệt và đồng bộ lại các workload tương ứng vào cluster.

\begin{figure}[H]
\centering
\reportimg{member1-report/06-argocd-pods-running.jpg}
\caption{Các pod ArgoCD trong namespace argocd đang ở trạng thái Running}
\end{figure}

Hai application chính là \texttt{yas-dev} và \texttt{yas-staging} được ArgoCD theo dõi. Ảnh chụp cho thấy cả hai application đạt trạng thái \texttt{Synced} và \texttt{Healthy}, xác nhận desired state trong Git đã khớp với trạng thái runtime của cluster tại thời điểm kiểm chứng.

\begin{figure}[H]
\centering
\reportimg{member1-report/07-argocd-applications-synced-healthy.jpg}
\caption{Application yas-dev và yas-staging được ArgoCD đồng bộ và đạt trạng thái Synced/Healthy}
\end{figure}

\subsection{Dịch vụ hạ tầng cho ứng dụng}

Các service hạ tầng được triển khai theo từng môi trường để hỗ trợ các workload YAS. Namespace \texttt{dev} chứa PostgreSQL, Kafka, Kafka Connect, Keycloak, Redis và Elasticsearch; đồng thời được dùng chung bởi luồng \texttt{developer-build}. Namespace \texttt{staging} có nhóm hạ tầng riêng để kiểm thử release độc lập với dữ liệu dev.

\begin{figure}[H]
\centering
\reportimg{member1-report/08-dev-infrastructure-services.jpg}
\caption{Các dịch vụ hạ tầng trong namespace dev đang Running và phục vụ môi trường dev/developer-build}
\end{figure}

\begin{figure}[H]
\centering
\reportimg{member1-report/09-staging-infrastructure-services.jpg}
\caption{Các dịch vụ hạ tầng trong namespace staging đang Running để phục vụ kiểm thử release}
\end{figure}

PostgreSQL lưu dữ liệu nghiệp vụ cho các service. Kafka và Kafka Connect hỗ trợ luồng đồng bộ dữ liệu bất đồng bộ như Debezium CDC. Keycloak cung cấp identity và authentication. Redis hỗ trợ cache, còn Elasticsearch phục vụ tìm kiếm sản phẩm. Nhóm hạ tầng này tạo nền tảng để các phần CI/CD, GitOps và Service Mesh hoạt động trên cùng một cluster.

\section{CI/CD Với Jenkins}

\subsection{Phạm vi pipeline}

Pipeline CI/CD phục vụ 16 workload ứng dụng trong scope hiện tại:

\begin{table}[H]
\centering
\caption{Danh sách workload trong scope pipeline}
\begin{tabular}{|p{0.32\textwidth}|p{0.22\textwidth}|p{0.34\textwidth}|}
\hline
\textbf{Workload} & \textbf{Source} & \textbf{Image} \\
\hline
\texttt{product}, \texttt{cart}, \texttt{order}, \texttt{customer}, \texttt{inventory}, \texttt{tax}, \texttt{payment}, \texttt{media}, \texttt{search}, \texttt{location}, \texttt{sampledata}, \texttt{storefront-bff}, \texttt{backoffice-bff} & cùng tên thư mục & \texttt{bingsu1103/<tên>:<tag>} \\
\hline
\texttt{storefront-ui} & \texttt{storefront/} & \texttt{bingsu1103/storefront:<tag>} \\
\texttt{backoffice-ui} & \texttt{backoffice/} & \texttt{bingsu1103/backoffice:<tag>} \\
\hline
\texttt{swagger-ui} & không build & \texttt{swaggerapi/swagger-ui} (public) \\
\texttt{kafka-connect} & không build & \texttt{quay.io/debezium/connect:2.4} (public) \\
\hline
\end{tabular}
\end{table}

Các service ngoài scope pipeline hiện tại gồm \texttt{promotion}, \texttt{rating}, \texttt{delivery}, \texttt{recommendation}, \texttt{webhook} và \texttt{payment-paypal}.

Pipeline được kích hoạt qua GitHub webhook hoặc Multibranch branch scan. Cơ chế monorepo detection cho phép chỉ build service có thay đổi, giảm thời gian chạy pipeline khi chỉ một phần nhỏ của repo bị thay đổi.

\subsection{Cấu hình Jenkins}

Jenkins sử dụng Multibranch Pipeline trỏ tới repo YAS và đọc \texttt{Jenkinsfile.ci} làm pipeline script. Agent chạy trên label \texttt{gcp-k8s-agent}, đảm bảo pipeline thực thi trên node có đủ công cụ build. Các credentials cần thiết gồm Docker Hub (\texttt{dockerhub-cred}), GitHub token (\texttt{github-pat}), kubeconfig (\texttt{gcp-kubeconfig}), SonarQube và Snyk token.

\begin{figure}[H]
\centering
\reportimg{member2-report/01-jenkins-dashboard.png}
\caption{Jenkins dashboard hiển thị các job CI/CD của đồ án}
\end{figure}

\begin{figure}[H]
\centering
\reportimg{member2-report/02-multibranch-config.png}
\caption{Multibranch Pipeline trỏ tới repo YAS và sử dụng \texttt{Jenkinsfile.ci}}
\end{figure}

\begin{figure}[H]
\centering
\reportimg{member2-report/03-jenkins-credentials.png}
\caption{Các credentials gồm Docker Hub, GitHub token, kubeconfig, SonarQube và Snyk}
\end{figure}

\subsection{Các stage trong CI Pipeline}

Pipeline CI/CD gồm các stage thực thi tuần tự:

\begin{enumerate}
\item \textbf{Pre-check}: kiểm tra môi trường build, xác nhận \texttt{java}, \texttt{mvn}, \texttt{gitleaks} và \texttt{snyk} sẵn sàng.
\item \textbf{Check Skip}: phát hiện file thay đổi bằng \texttt{git diff} so với \texttt{origin/main}. Nếu chỉ có file tài liệu (\texttt{docs/}, \texttt{*.md}), pipeline bỏ qua các stage CI/CD. Trên main branch, khi \texttt{merge-base == HEAD}, pipeline sử dụng fallback \texttt{HEAD\textasciitilde1..HEAD} để tránh diff rỗng.
\item \textbf{Secret Scanning}: quét secret bằng \texttt{gitleaks} trên phạm vi commit thay đổi.
\item \textbf{Monorepo Execution}: với mỗi backend service, kiểm tra thư mục có thay đổi hay không. Nếu có, chạy \texttt{mvn test} và \texttt{mvn package}. Nếu có release tag hoặc thay đổi \texttt{common-library}, build toàn bộ service.
\item \textbf{Code Quality / Quality Gate}: gửi kết quả phân tích lên SonarQube. Nếu SonarQube không khả dụng, pipeline ghi cảnh báo và tiếp tục các stage CD.
\item \textbf{Coverage Report}: publish báo cáo JaCoCo cho từng service có thay đổi. Ngưỡng instruction coverage tối thiểu là 70\%.
\item \textbf{Dependency Scan}: quét dependency bằng Snyk, lưu báo cáo và monitor project.
\item \textbf{Docker Build \& Push}: build image và push lên Docker Hub (chi tiết ở mục tiếp theo).
\item \textbf{Update GitOps}: cập nhật image tag trong repo \texttt{gitops-manifest-k8s} (chi tiết ở Mục~\ref{sec:gitops}).
\end{enumerate}

\begin{figure}[H]
\centering
\reportimg{member2-report/04-pipeline-stages.png}
\caption{Jenkins pipeline hiển thị các stage CI/CD chính}
\end{figure}

\begin{figure}[H]
\centering
\reportimg{member2-report/05-pipeline-console-output.png}
\caption{Console output cho thấy service được detect, build và push image}
\end{figure}

\subsection{Docker Build và Push}

Image được tag theo short commit SHA (7 ký tự đầu). Trên branch \texttt{main}, image được tag thêm \texttt{latest}. Khi có release tag dạng \texttt{vX.Y.Z} tại HEAD, image được tag thêm theo tên tag đó. Tất cả image được push lên Docker Hub với prefix \texttt{bingsu1103/}.

Với UI workload, source code nằm trong thư mục \texttt{storefront/} và \texttt{backoffice/} nhưng image được đặt tên là \texttt{bingsu1103/storefront} và \texttt{bingsu1103/backoffice}. Pipeline phát hiện thay đổi dựa trên đường dẫn source, không dựa trên tên chart Kubernetes.

\begin{figure}[H]
\centering
\reportimg{member2-report/06-dockerhub-product-tags.jpg}
\caption{Docker Hub hiển thị image tag của service \texttt{product}}
\end{figure}

\begin{figure}[H]
\centering
\reportimg{member2-report/07-dockerhub-location-tags.jpg}
\caption{Docker Hub hiển thị image tag của service \texttt{location}}
\end{figure}

\begin{figure}[H]
\centering
\reportimg{member2-report/08-dockerhub-ui-tags.jpg}
\caption{Docker Hub hiển thị image \texttt{storefront}/\texttt{backoffice} dùng cho UI workload}
\end{figure}

\subsection{Developer Build}

Ngoài pipeline CI/CD chính, nhóm cấu hình job \texttt{developer-build} cho phép build và triển khai tạm thời từ branch phát triển. Job này nhận tham số là tên branch cho từng service trong scope. Pipeline resolve image tag từ commit ID trên remote branch, sau đó deploy vào namespace \texttt{developer-build} bằng Kustomize và \texttt{kubectl apply}.

Namespace \texttt{developer-build} sử dụng NodePort Service để developer có thể truy cập service qua \texttt{<WORKER\_IP>:<NODE\_PORT>}. Deployment trong namespace này được dọn dẹp sau khi kiểm thử hoàn tất để tiết kiệm tài nguyên cluster.

\begin{figure}[H]
\centering
\reportimg{member2-report/11-developer-build-parameters.png}
\caption{Job \texttt{developer\_build} cho phép chọn branch cho từng workload trong scope}
\end{figure}

\begin{figure}[H]
\centering
\reportimg{member2-report/12-developer-build-nodeport.png}
\caption[Kết quả deploy và rollout trong namespace developer-build]{Console output hiển thị kết quả deploy và trạng thái rollout trong namespace \texttt{developer-build}}
\end{figure}

\subsection{Cleanup Job}

Job \texttt{cleanup} xóa toàn bộ Deployment, Service, ConfigMap, ReplicaSet và Pod trong namespace \texttt{developer-build}. Namespace và Secret được giữ lại để tái sử dụng cho lần deploy tiếp theo. Sau khi xóa, job chạy \texttt{kubectl get all} để xác nhận không còn workload nào trong namespace.

\begin{figure}[H]
\centering
\reportimg{member2-report/13-cleanup-console.png}
\caption{Cleanup job xóa workload trong namespace \texttt{developer-build}}
\end{figure}

\begin{figure}[H]
\centering
\reportimg{member2-report/14-developer-build-after-cleanup.png}
\caption{Namespace \texttt{developer-build} sau cleanup, không còn Deployment hoặc Pod}
\end{figure}

\subsection{Sự cố và cách xử lý}

Trong quá trình phát triển pipeline, nhóm gặp một số sự cố cần xử lý:

\begin{itemize}
\item \textbf{Diff rỗng trên main branch}: khi \texttt{merge-base} trùng với \texttt{HEAD}, \texttt{git diff} không trả kết quả. Pipeline đã bổ sung fallback \texttt{HEAD\textasciitilde1..HEAD} để phát hiện thay đổi trong commit mới nhất.
\item \textbf{UI source mapping}: \texttt{storefront-ui} và \texttt{backoffice-ui} build từ thư mục \texttt{storefront/} và \texttt{backoffice/}, không phải từ tên chart Kubernetes. Pipeline cần map đúng source path để detect thay đổi.
\item \textbf{Image public}: \texttt{swagger-ui} dùng image \texttt{swaggerapi/swagger-ui} và \texttt{kafka-connect} dùng image \texttt{quay.io/debezium/connect:2.4}. Hai workload này không cần build từ repo YAS.
\end{itemize}

\section{GitOps Và Kustomize}

Repo \texttt{gitops-manifest-k8s} là nguồn cấu hình triển khai cho Kubernetes. Repo này tổ chức manifest theo mô hình \texttt{base} và \texttt{overlays}. \texttt{base} định nghĩa Deployment, Service, ServiceAccount và cấu hình chung của từng workload. Các overlay \texttt{dev}, \texttt{staging} và \texttt{developer-build} điều chỉnh image tag, namespace, replica và cấu hình môi trường.

ArgoCD theo dõi repo GitOps và đồng bộ trạng thái mong muốn vào cluster. Khi Jenkins cập nhật image tag trong repo manifest, ArgoCD phát hiện thay đổi và rollout workload tương ứng. Cơ chế này bảo đảm quy trình CD có lịch sử rõ ràng, có thể truy vết qua commit GitOps và trạng thái ArgoCD.

Minh chứng cho phần này gồm cấu trúc repo GitOps, kết quả Kustomize render, trạng thái ArgoCD sync và workload trong từng namespace.

\section{Service Mesh Với Istio, Security Policy Và Kiali}

\subsection{Mục tiêu và phạm vi}

Phần Service Mesh tập trung vào cấu hình và kiểm chứng luồng giao tiếp giữa các service của hệ thống YAS bằng Istio. Phạm vi chính gồm ingress routing, sidecar injection, mTLS, DestinationRule, retry policy, AuthorizationPolicy, kiểm thử allow/deny và trực quan hóa bằng Kiali. Môi trường \texttt{dev} được dùng để kiểm chứng chi tiết các chính sách mesh. Môi trường \texttt{staging} được bổ sung overlay Istio riêng với mTLS, DestinationRule, retry policy và AuthorizationPolicy tương ứng để đáp ứng yêu cầu triển khai ở môi trường tiền sản xuất.

\subsection{Cấu hình Istio}

Cấu hình Istio được quản lý theo hướng GitOps trong repo manifest. Gateway và VirtualService ở lớp base định tuyến request từ ingress gateway tới Storefront, Swagger UI và các BFF/API tương ứng. Overlay \texttt{dev} và \texttt{staging} bổ sung các resource bảo mật và điều phối traffic, gồm \texttt{PeerAuthentication}, \texttt{DestinationRule}, \texttt{VirtualService} retry và \texttt{AuthorizationPolicy}. Các file staging được tách riêng tại \texttt{environments/staging/istio} để host và ServiceAccount principal trỏ đúng namespace \texttt{staging}.

\begin{figure}[H]
\centering
\reportimg{member4-report/01-kiali-istio-config-overview.png}
\caption{Kiali Istio Config hiển thị các resource đại diện trong namespace \texttt{dev}}
\end{figure}

Hình trên cho thấy các resource Istio chính tồn tại trong namespace \texttt{dev} và có trạng thái hợp lệ. Các nhóm cấu hình đại diện gồm access policy, mTLS, routing và retry. Đây là bằng chứng trực quan cho việc cấu hình Istio đã được đồng bộ vào cluster.

\subsection{Sidecar injection và trạng thái Istio system}

Istio system, Kiali, Prometheus và Grafana được triển khai trong namespace \texttt{istio-system}. Namespace \texttt{dev}, \texttt{staging} và \texttt{developer-build} được gắn label \texttt{istio-injection=enabled} để Istio tự động inject sidecar vào pod ứng dụng.

\begin{figure}[H]
\centering
\reportimg{member4-report/02-istio-system-pods-services.png}
\caption{Các thành phần Istio, Kiali, Prometheus và Grafana đang chạy trong \texttt{istio-system}}
\end{figure}

\begin{figure}[H]
\centering
\reportimg{member4-report/03-namespace-sidecar-injection-labels.png}
\caption{Các namespace triển khai có label \texttt{istio-injection=enabled}}
\end{figure}

Kết quả kiểm tra pod trong namespace \texttt{dev} cho thấy mỗi pod ứng dụng có container nghiệp vụ và container \texttt{istio-proxy}. Trạng thái \texttt{2/2 Running} chứng minh workload và sidecar đều sẵn sàng.

\begin{figure}[H]
\centering
\reportimg{member4-report/04-dev-pods-sidecar-containers.png}
\caption{Pod trong \texttt{dev} có container ứng dụng và sidecar \texttt{istio-proxy}}
\end{figure}

\begin{figure}[H]
\centering
\reportimg{member4-report/05-dev-pods-running-2of2.png}
\caption{Các pod ứng dụng trong \texttt{dev} ở trạng thái \texttt{2/2 Running}}
\end{figure}

\subsection{mTLS và DestinationRule}

Namespace \texttt{dev} sử dụng \texttt{PeerAuthentication} ở chế độ \texttt{STRICT}. Cấu hình này yêu cầu traffic service-to-service trong mesh phải sử dụng mTLS. Các \texttt{DestinationRule} tương ứng cấu hình \texttt{ISTIO\_MUTUAL} cho các service có sidecar, giúp client trong mesh sử dụng TLS do Istio quản lý khi gọi nội bộ.

\begin{figure}[H]
\centering
\reportimg{member4-report/06-mtls-peer-authentication-destinationrules.png}
\caption{PeerAuthentication STRICT và DestinationRule ISTIO\_MUTUAL trong namespace \texttt{dev}}
\end{figure}

Cách cấu hình này giới hạn phạm vi áp dụng vào các service trong mesh, tránh ảnh hưởng tới thành phần hạ tầng hoặc workload không có sidecar.

\subsection{Retry policy}

Retry policy được triển khai bằng \texttt{VirtualService}. Các service quan trọng trong luồng nghiệp vụ như \texttt{product}, \texttt{cart}, \texttt{order}, \texttt{tax}, \texttt{payment} và \texttt{inventory} có cấu hình retry cho lỗi tạm thời. Ví dụ minh chứng là \texttt{tax-retry} với số lần thử lại và thời gian timeout trên từng lần thử.

\begin{figure}[H]
\centering
\reportimg{member4-report/07-virtualservice-tax-retry.png}
\caption{VirtualService \texttt{tax-retry} cấu hình retry cho lỗi tạm thời}
\end{figure}

Retry policy giúp tăng độ ổn định của service-to-service traffic trong trường hợp lỗi mạng ngắn hạn hoặc backend phản hồi không ổn định trong thời gian ngắn.

\subsection{AuthorizationPolicy và kiểm thử allow/deny}

AuthorizationPolicy được dùng để giới hạn quyền truy cập giữa các service dựa trên ServiceAccount. Cách tiếp cận này phù hợp với microservices vì mỗi workload có ServiceAccount riêng, từ đó chính sách có thể định nghĩa rõ service nào được gọi service nào.

Nhóm kiểm thử hai trường hợp đại diện. Trường hợp hợp lệ sử dụng ServiceAccount \texttt{storefront-bff} gọi endpoint health của \texttt{product} và nhận HTTP \texttt{200}. Trường hợp không hợp lệ sử dụng ServiceAccount \texttt{search} gọi \texttt{payment} và bị Istio RBAC chặn với HTTP \texttt{403}.

\begin{figure}[H]
\centering
\reportimg{member4-report/08-authz-allow-deny-curl-test.png}
\caption{Kiểm thử AuthorizationPolicy: request hợp lệ trả \texttt{200}, request không hợp lệ bị chặn với \texttt{403 RBAC}}
\end{figure}

Kết quả này chứng minh policy không chỉ tồn tại ở manifest mà đã có hiệu lực ở runtime.

\subsection{Kiali visualization}

Kiali được sử dụng để quan sát topology, traffic và trạng thái bảo mật trong service mesh. Overview cho biết namespace nằm trong mesh và graph thể hiện quan hệ gọi giữa các service. Trong Service Graph, biểu tượng khóa trên các cạnh kết nối là minh chứng trực quan cho traffic được bảo vệ bằng mTLS.

\begin{figure}[H]
\centering
\reportimg{member4-report/09-kiali-overview-namespaces.png}
\caption{Kiali Overview hiển thị các namespace trong mesh}
\end{figure}

\begin{figure}[H]
\centering
\reportimg{member4-report/10-kiali-service-graph-mtls.png}
\caption{Kiali Service Graph thể hiện topology, traffic và mTLS trong namespace \texttt{dev}}
\end{figure}

Kiali bổ sung góc nhìn vận hành cho phần cấu hình YAML. Thông qua graph, nhóm có thể kiểm tra service nào đang phát sinh traffic, kết nối nào có mTLS và workload nào cần điều chỉnh policy.

\subsection{Kiểm chứng Istio ingress}

Istio ingress gateway được kiểm chứng bằng cách port-forward service \texttt{istio-ingressgateway} và gửi request với Host header tương ứng. Các endpoint đại diện gồm Storefront, Swagger UI, Product categories và Payment providers đều trả HTTP \texttt{200}.

\begin{figure}[H]
\centering
\reportimg{member4-report/11-ingress-curl-checks.png}
\caption{Kiểm chứng ingress: các endpoint chính đều trả HTTP \texttt{200}}
\end{figure}

Kết quả này xác nhận Gateway và VirtualService đã định tuyến đúng request từ ingress vào các service trong hệ thống.

\subsection{Kết luận phần Service Mesh}

Phần Service Mesh đã hoàn thành các yêu cầu chính: sidecar injection, mTLS, DestinationRule, retry policy, AuthorizationPolicy, kiểm thử allow/deny, Kiali visualization và ingress routing. Các minh chứng bằng hình ảnh cho thấy cấu hình trong \texttt{dev} đã được apply và có hiệu lực ở runtime. Overlay \texttt{staging} đã được tách riêng trong repo GitOps để áp dụng cùng nhóm chính sách mesh cho môi trường tiền triển khai mà không phụ thuộc trực tiếp vào cấu hình namespace \texttt{dev}.

\section{Observability}

Nhóm đã triển khai bổ sung Observability trong namespace \texttt{observability}. Stack gồm Prometheus cho metrics, Grafana cho dashboard, Loki và Promtail cho log, Tempo cho distributed tracing và OpenTelemetry Collector cho việc tiếp nhận, xử lý, chuyển tiếp telemetry.

Phần này là khung để bổ sung hình ảnh sau khi nhóm chụp UI Grafana. Các minh chứng cần ưu tiên gồm: Helm releases trong namespace \texttt{observability}, trạng thái pod/service, Prometheus targets hoặc query, Grafana datasources, JVM dashboard, Loki logs và Tempo traces. Nội dung chi tiết được chuẩn bị trong file \texttt{docs/observability-report.md}.

\section{Kết Luận}

Đồ án đã xây dựng được quy trình triển khai theo hướng CI/CD và GitOps cho hệ thống YAS. Jenkins đảm nhiệm build/test/push image và cập nhật manifest. ArgoCD đồng bộ trạng thái triển khai từ repo GitOps vào cluster K3s. Istio được sử dụng để bổ sung Service Mesh cho môi trường \texttt{dev} và \texttt{staging}, bao gồm mTLS, retry, AuthorizationPolicy và ingress routing. Kiali và Observability stack cung cấp minh chứng trực quan cho trạng thái vận hành của hệ thống.

Kết quả hiện tại cho thấy nhóm đã triển khai được pipeline và hạ tầng cần thiết để vận hành các workload chính trong scope. Các bước tiếp theo khi hoàn thiện báo cáo là bổ sung đầy đủ hình ảnh minh chứng cho phần hạ tầng, CI/CD, GitOps và Observability, sau đó rà soát lại số hiệu hình, caption và trạng thái cluster tại thời điểm nộp.


\end{document}
